\lecture{15}{Tue 25 Mar 2025 14:04}{More Schrodinger Equation}
Ok they are actually partial derivatives because $\psi (x,t)$ since $\psi $ has time dependence. We usually write it as a total derivative though, as we usually write the state as $\ket{\psi ,t} $ for some reason.

To show
\[
	\left\langle P \right\rangle =m \frac{\mathrm{d}}{\mathrm{d}t} \left\langle X \right\rangle 
,\]
we write
\[
	\frac{\mathrm{d}}{\mathrm{d}t} \bra{\psi } X \ket{\psi } =\frac{\mathrm{d}}{\mathrm{d}t} \bra{\psi } \int x \ket{x} \bra{x}  \,\mathrm{d} x \ket{\psi } =\frac{\mathrm{d}}{\mathrm{d}t} \int \braket{\psi }{x}x \braket{x}{\psi }  \,\mathrm{d} x=\frac{\mathrm{d}}{\mathrm{d}t} \int \psi ^*(x) x \psi (x) \,\mathrm{d} x
.\]
We write this as
\[
	\frac{\mathrm{d}}{\mathrm{d}t} \bra{\psi ,t} X \ket{\psi ,t} 
.\]
The time dependence of states comes from the Schrodinger equation
\[
	i \hbar \frac{\mathrm{d}}{\mathrm{d}t} \ket{\psi ,t} =H \ket{\psi ,t} =\frac{P^2}{2m} \ket{\psi ,t} +V(x)\ket{\psi ,t} 
.\]
Taking the Hermitian conjugate gives us
\[
	i \hbar \frac{\mathrm{d}}{\mathrm{d}t} \bra{\psi ,t} =-\bra{\psi ,t} H
,\]
because $H$ is Hermitian. Using our above equation, we have
\[
	\frac{\mathrm{d}}{\mathrm{d}t} \bra{\psi ,t} X\ket{\psi ,t} =\left( \frac{\mathrm{d}}{\mathrm{d}t} \bra{\psi ,t}  \right) X\ket{\psi ,t} +\bra{\psi ,t} X\left( \frac{\mathrm{d}}{\mathrm{d}t} \ket{\psi ,t}  \right) 
.\]
Plugging this in, we obtain
\[
\frac{\mathrm{d}}{\mathrm{d}t} \left\langle X \right\rangle =\frac{1}{i\hbar } \left( -\bra{\psi ,t} HX\ket{\psi ,t} +\bra{\psi ,t} XH\ket{\psi ,t}  \right) =-\frac{1}{i\hbar } \bra{\psi ,t} \left( -HX+XH \right) \ket{\psi ,t} =\frac{1}{i\hbar } \left\langle \left[ X,H \right]  \right\rangle 
.\]
This generalizes to any \emph{time-independent} operator $X$, not just position. This shows us that
\[
	\frac{\mathrm{d}}{\mathrm{d}t} \left\langle A \right\rangle =\frac{1}{i\hbar } \left\langle \left[ A,H \right] \right\rangle 
.\]
Note that $\left[ A,BC \right] =\left[ A,B \right] C+B\left[ A,C \right] $.
\[
	\left[ X,H \right] =\left[ X,\frac{p^2}{2m} +V(x) \right] =\left[ X,p^2/2m \right] +\left[ X,V\left( x \right)  \right] 
.\]
Recall that $\left[ X,p\right]=i\hbar  $. We have
\[
	\frac{1}{2m} \left(\left[ X,p \right] p+p\left[ X,p \right] \right)=\frac{1}{2m} \left(i\hbar p+i\hbar p\right)=\frac{i\hbar }{m} p
.\]
Additionally, $\left[ X,V(x) \right] =0$, since they share eigenstates. We can plug this back in to obtain
\[
	\frac{\mathrm{d}}{\mathrm{d}t} \left\langle X \right\rangle =\frac{1}{i\hbar } \left\langle \left[ X,H \right]  \right\rangle =\frac{1}{i\hbar } \frac{i\hbar }{m} \left\langle p \right\rangle =\frac{1}{m} \left\langle p \right\rangle 
.\]
Therefore, $m\frac{\mathrm{d}X}{\mathrm{d}t} =P$. This satisfies the classical relation for position and momentum, so we have that the momentum operator behaves as expected.

\section{Solving Schrodinger's Equation}
We will take various potentials $V(x)$, and solve the Schrodinger equation in that potential. The simplest one is in fact the infinite square well:
\[
	V(x)=
	\begin{cases}
		0,&0<x<L\\
		\infty ,&\text{otherwise} 
	\end{cases}
.\]
We want to find all eigenvals/vecs of $H$ in this potential:
\[
	H=\frac{p^2}{2m}+ V(x)
.\]
We then have the eigenvalue problem
\[
	H\ket{\psi } \cong \left( -\frac{\hbar ^2}{2m} \frac{\mathrm{d}^2}{\mathrm{d}x^2} +V(x) \right) \psi (x)=E\ket{\psi } 
.\]
Since $V$ goes to infinity outside of $(0,L)$, the wavefunction must be 0 in this interval, since the particle cannot go there. We can then impose boundary conditions on our equation. We will start by solving it in the region where $V(x)=0$:
\[
	-\frac{\hbar ^2}{2m} \frac{\mathrm{d}^2\psi }{\mathrm{d}x^2} =E\psi (x)\implies \frac{\mathrm{d}^2\psi }{\mathrm{d}x^2} =-\frac{2mE}{\hbar ^2} \psi (x)
.\]
Set $k=\sqrt{\frac{2mE}{\hbar ^2} } $. We then have
\[
	\psi (x)=\exp \left( \pm ikx \right) 
.\]
Our general solution is thus
\[
	\psi (x)=A'e^{ikx} +B'e^{-ikx} 
.\]
To identify $E$ from $k$, we impose the boundary conditions. We need to find the values of $k$ such that we oscillate precisely an integer number of times within the well, since otherwise it would be nonzero on the boundary (the wavefunction needs to be continuous).

Our boundary conditions are that $\psi (0)=\psi (L)=0$. Plugging this in, we obtain
\[
	0=A'e^{0} +B'e^{0} =A'+B'
.\]
Therefore, $A'=-B'$. Plugging in $L$, we have
\[
	0=A'e^{ikL} +B'e^{-ikL} =A'e^{ikL} -A'e^{-ikL} 
.\]
This forces $e^{ikL} -e^{-ikL} =0$, and thus
\[
	\cos kL+i \sin kL-\cos (-kL)-i \sin (-kL)=2i \sin kL
.\]
This quantizes $k$ immediately, as $\psi (x)\neq 0$ since the particle exists, and thus
\[
	kL\in \pi\mathbb{Z} \implies \sqrt{\frac{2mE}{\hbar ^2} } L=\pi n\implies E=\frac{n^2}{L^2} \frac{\hbar ^2}{2m}\pi^2 =\frac{n^2\hbar ^2\pi^2 }{2Lm}  
\]
for any $n\in\mathbb{Z}^+ $. We cannot have $n<0$ for obvious reasons. Therefore, after normalization we have
\[
	\psi (x)=A \sin \left( \frac{n\pi x}{L}  \right) =A \sin \left( k_nx \right) ,\qquad k_n\coloneqq \frac{n\pi }{L} 
.\]
The sin instead of the exponential comes from the decomposition into $\sin $ we did earlier, but more general. The normalization $A'\to A$ is given by
\[
	1=\braket{\psi }{\psi } =\int_{\mathbb{R} }  \left| \psi (x) \right| ^2 \,\mathrm{d} x\implies \braket{\psi _n}{\psi _n} =\int_{0}^{L} \left| A \sin (k_nx) \right| ^2 \,\mathrm{d} x=\left| A \right| ^2\int_{0}^{L} \sin ^2(k_nx) \,\mathrm{d} x
.\]
This integral evaluates to $L/2$, giving the normalization factor for $A$. We then have $A=\sqrt{2/L} $, and each $\ket{\psi _n} $ has norm 1:
\[
	\psi _n(x)=\sqrt{\frac{2}{L} } \sin \left( k_nx \right)
.\]
These are the solutions to the Schrodinger equation. Since the energy levels grow like $E_n\propto n^2$, we have $E_n=n^2E_1$. Since these functions are all symmetric, the probability in the left/right halves of the box are 1/2.
